\cotefile{84}
\scanpage{97}

\diagram{\begin{tikzcd}[ampersand replacement=\&,column sep=large]
0\arrow[r] \& L\arrow[r,dash]\arrow[d,"{-\mathrm{id}}"] \& E\arrow[r]\arrow[d,"g","?"'] \& \underline{\mathcal O}\arrow[r]\arrow[d,"\mathrm{id}"] \& 0\\
0\arrow[r,dash] \& L\arrow[r] \& A\otimes L\arrow[r,"{\varphi\otimes\mathrm{id}_{L}}"] \& \underline{\mathcal O}\arrow[r] \& 0
\end{tikzcd}}

\margin{Attention au signe $\longrightarrow$}

Prouvons d'abord

\[
\varphi\otimes\mathrm{id}_{L}\bigl(\underbrace{g(x)}_{\substack{\rotatebox{90}{$\in$}\\ \deleted{\unsure{\mathrm{H}|\mathrm{Hom}}}A\otimes L=\mathrm{Hom}(E,L)}}\bigr)=\psi(x)
\qquad \forall x\in\Gamma(E)
\]

Cela signifie que $\forall x\in E$,

\[
g(x)|L=(u\longmapsto\psi(x)u)
\]

\deleted{\illegible{1}} i.e. on a

\[
\underset{\substack{\Vert\\ \psi(x).u-\underset{\substack{\Vert\\0}}{\psi(u)}x}}{g(x)(u)}
=\psi(x)u\qquad\text{pour tout }u\in\Gamma L\subset\Gamma E
\]

c'est OK.

Prouvons ensuite que pour $x\deleted{\unsure{\in E}}\in\Gamma L$, on a $g(x)\unsure{:}E\longrightarrow L$ est \uline{nul} sur $L$, et par passage au quo. induit $\lambda\longmapsto-\lambda x$ de $\mathcal O$ en \unsure{$x$|$L$},

\margin{signe $\nearrow$}

\[
\text{i.e.}\qquad g(x)(y)=-\psi(y)x,
\]

mais c'est aussi OK.

Ceci vu, l'isom. d'ext. $g$ définit

\[
\mathrm{Quad}(E,L^{\otimes2})\simeq\mathrm{Quad}(A,\mathcal O)
\]

\deleted{$\unsure{\mathrm{Sym}^{\! \unsure{2|4}}|\mathrm{Quad}}( \check E)\otimes L^{\otimes2}\qquad \unsure{\mathrm{Sym}^{2}}$}.

% slanted
\uline{dis que \inserted{lorsque $E$ provient d'un \unsure{rev.} du \unsure{1\textsuperscript{er}} type, d'où $N$ sur $A$, $Q$ sur $E$, alors} par cet} iso, \uline{$N$ et} \uline{$Q$ se} \uline{correspondent}. Pour le prouver, OPS $E$ scindé en

\[
E\simeq L\oplus\underline{\mathcal O}.x_0
\]

\[
\text{d'où}\qquad A\simeq\underline{\mathcal O}\oplus L^\vee \qquad
\text{i.e.}\quad A\otimes L\simeq L\oplus
\underbrace{L\otimes L^\vee}_{\underline{\mathcal O}_X}
\]

% archive sheet number: 96

